\documentclass[12pt]{article}
\usepackage{amsmath, amssymb, amsthm}
\usepackage{geometry}
\geometry{margin=1in}
 
\newtheorem*{problem}{Problem}
\newtheorem*{solution}{Solution}
 
\title{Linear Algebra Done Right (4th Ed.) -- Axler\\[0.5em]
       \large Matrix of the Differentiation Map}
\author{}
\date{}
 
\begin{document}
\maketitle
 
\begin{problem}
Suppose $D \in \mathcal{L}(\mathcal{P}_3(\mathbb{R}),\, \mathcal{P}_2(\mathbb{R}))$ is the
differentiation map defined by $Dp = p'$. Find a basis of $\mathcal{P}_3(\mathbb{R})$ and a
basis of $\mathcal{P}_2(\mathbb{R})$ such that the matrix of $D$ with respect to these bases
is
\[
  \begin{pmatrix}
    1 & 0 & 0 & 0 \\
    0 & 1 & 0 & 0 \\
    0 & 0 & 1 & 0
  \end{pmatrix}.
\]
\end{problem}
 
\begin{solution}
Recall that if $T \in \mathcal{L}(V, W)$, and $v_1, \ldots, v_n$ and $w_1, \ldots, w_m$ are
bases of $V$ and $W$ respectively, then the $k$-th column of $\mathcal{M}(T)$ records the
coordinates of $Tv_k$ in terms of the $w_j$ basis. In particular, the given matrix requires
\[
  Dv_k = w_k \quad \text{for } k = 1, 2, 3, \qquad \text{and} \qquad Dv_4 = 0.
\]
 
\medskip
\noindent\textbf{Choosing a basis for $\mathcal{P}_3(\mathbb{R})$.}
Since $\dim \mathcal{P}_3(\mathbb{R}) = 4$, we need four linearly independent polynomials.
We take the standard monomial basis:
\[
  v_1 = x^3,\quad v_2 = x^2,\quad v_3 = x,\quad v_4 = 1.
\]
Computing the derivative of each basis vector gives
\[
  Dv_1 = 3x^2, \qquad Dv_2 = 2x, \qquad Dv_3 = 1, \qquad Dv_4 = 0.
\]
 
\medskip
\noindent\textbf{Choosing a basis for $\mathcal{P}_2(\mathbb{R})$.}
We need the basis vectors $w_1, w_2, w_3$ of $\mathcal{P}_2(\mathbb{R})$ to satisfy
$w_k = Dv_k$. The computations above dictate the choice:
\[
  w_1 = 3x^2, \qquad w_2 = 2x, \qquad w_3 = 1.
\]
Since $\dim \mathcal{P}_2(\mathbb{R}) = 3$ and these three polynomials are clearly linearly
independent (they have distinct degrees), they form a basis of $\mathcal{P}_2(\mathbb{R})$.
 
\medskip
\noindent\textbf{Verification.}
With these choices, the action of $D$ on each basis vector of $\mathcal{P}_3(\mathbb{R})$ is:
\[
  Dv_1 = 3x^2 = 1\cdot w_1 + 0\cdot w_2 + 0\cdot w_3,
\]
\[
  Dv_2 = 2x   = 0\cdot w_1 + 1\cdot w_2 + 0\cdot w_3,
\]
\[
  Dv_3 = 1    = 0\cdot w_1 + 0\cdot w_2 + 1\cdot w_3,
\]
\[
  Dv_4 = 0    = 0\cdot w_1 + 0\cdot w_2 + 0\cdot w_3.
\]
Reading off the coordinate columns yields exactly the desired matrix:
\[
  \mathcal{M}(D) =
  \begin{pmatrix}
    1 & 0 & 0 & 0 \\
    0 & 1 & 0 & 0 \\
    0 & 0 & 1 & 0
  \end{pmatrix}.
\]
 
\medskip
\noindent\textbf{Comparison with Example 3.33.}
In Example 3.33, the standard monomial bases $1, x, x^2, x^3$ and $1, x, x^2$ are used for
both spaces. Under those bases, the matrix of $D$ contains the scalar coefficients $1, 2, 3$
along its super-diagonal, reflecting the factors produced by differentiation. Here, by
choosing $w_k = Dv_k$ for each $k$, we \emph{absorb} those scalar factors directly into the
basis vectors of $\mathcal{P}_2(\mathbb{R})$, yielding the cleaner identity-like matrix above.
\end{solution}
 
\end{document}
